\subsubsection{Running Example}\label{sec:approach-example}

Consider a service desk observed over consecutive feature windows $W_{\mathrm a},\ldots,W_{\mathrm f}$. In each window $W_{\mathrm x}$, a focal case $\mathrm x$ and its companion case overlap in time. Each has a Receive event followed by a Resolve event, and both events fall within that window. The Frontline team receives every case and resolves every companion case after a gap of $g$. It also resolves the focal cases $\mathrm a$, $\mathrm b$, and $\mathrm c$ after the same gap. The Specialist team resolves the focal cases $\mathrm d$, $\mathrm e$, and $\mathrm f$ after a longer gap $G>g$.

\begin{Example}[Intra-case features for a case prefix]\label{ex:intra-case}
Compare the Receive and Resolve prefixes of cases $\mathrm c$ and $\mathrm d$. At Receive, the current-activity indicator for Receive is on and the prefix contains no directly-follows pair. At Resolve, Receive is the immediately preceding activity and the prefix contains the Receive-to-Resolve pair once. \autoref{tab:example-intra-case} shows the same feature values at the corresponding positions of both cases. \autoref{fig:example-intra-case} illustrates the prefixes used to build the rows.
\end{Example}

\begin{table*}[!t]
  \caption{Selected intra-case features for the Receive and Resolve events of cases $\mathrm c$ and $\mathrm d$. On and off denote true and false for an indicator.}
  \label{tab:example-intra-case}
  \centering
  \small
  \begin{tabularx}{\linewidth}{@{}ll*{3}{>{\centering\arraybackslash}X}@{}}
    \toprule
    \textbf{Case} & \textbf{Current event} & \textbf{Receive now} &
      \textbf{Receive immediately before} & \textbf{Receive-to-Resolve count} \\
    \midrule
    $\mathrm c$ & Receive & On & Off & Zero \\
    $\mathrm c$ & Resolve & Off & On & One \\
    $\mathrm d$ & Receive & On & Off & Zero \\
    $\mathrm d$ & Resolve & Off & On & One \\
    \bottomrule
  \end{tabularx}
\end{table*}

\begin{figure*}[!t]
  \centering
  \begin{tikzpicture}[
  kairo/base, x=1cm, y=1cm,
  activity/.style={
    rounded corners=2pt, minimum width=1.58cm, minimum height=0.37cm,
    draw=KairoLine, fill=white, inner sep=2pt
  },
  current/.style={activity, draw=KairoBlue, fill=KairoBlueTint, text=KairoBlue},
  follows/.style={-{Stealth[length=4pt]}, draw=KairoMuted, line width=0.45pt}
]
  \path[use as bounding box] (0,0) rectangle (13.08,-1.56);
  \path[fill=KairoPaper, rounded corners=3pt] (1.69,-1.55) rectangle (7.23,-0.40);
  \path[fill=KairoPaper, rounded corners=3pt] (7.49,-1.55) rectangle (13.03,-0.40);
  \node[anchor=west, font=\footnotesize\bfseries] at (0,-0.16) {Case};
  \node[font=\footnotesize\bfseries] at (4.46,-0.16) {Prefix at Receive};
  \node[font=\footnotesize\bfseries] at (10.26,-0.16) {Prefix at Resolve};
  \foreach \caseid/\rowy in {c/-0.72,d/-1.23} {
    \node[anchor=west] at (0,\rowy) {$\mathrm{\caseid}$};
    \node[current] at (4.46,\rowy) {Receive};
    \node[activity, text=KairoMuted] (receive-\caseid) at (9.29,\rowy) {Receive};
    \node[current] (resolve-\caseid) at (11.23,\rowy) {Resolve};
    \draw[follows] (receive-\caseid.east) -- (resolve-\caseid.west);
  }
\end{tikzpicture}

  \caption{Cases $\mathrm c$ and $\mathrm d$ have matching activity prefixes at each event. The shaded activity is the current event in each prefix.}
  \label{fig:example-intra-case}
\end{figure*}

\begin{Example}[Resource features for a handover]\label{ex:resource}
Between $W_{\mathrm c}$ and $W_{\mathrm d}$, responsibility for resolving the focal case moves from Frontline to Specialist. Frontline still resolves the companion case, so Specialist handles half of the Resolve events in $W_{\mathrm d}$. The only handover in that window goes from Frontline to Specialist, with a gap of $G$. \autoref{tab:example-resource} shows how the assignment and handover features change. \autoref{fig:example-resource} shows who performs each focal event.
\end{Example}

\begin{table*}[!t]
  \caption{Selected resource features in consecutive feature windows. The handover share and mean incoming gap are zero in $W_{\mathrm c}$, which contains no handover.}
  \label{tab:example-resource}
  \centering
  \small
  \begin{tabularx}{\linewidth}{@{}X*{2}{>{\centering\arraybackslash}p{0.17\linewidth}}@{}}
    \toprule
    \textbf{Feature} & $\boldsymbol{W_{\mathrm c}}$ & $\boldsymbol{W_{\mathrm d}}$ \\
    \midrule
    Events by Frontline & Four & Three \\
    Events by Specialist & Zero & One \\
    Share of Resolve performed by Specialist & Zero & Half \\
    Share of Frontline's handovers received by Specialist & Zero & Full \\
    Mean incoming handover gap for Specialist & Zero & $G$ \\
    \bottomrule
  \end{tabularx}
\end{table*}

\begin{figure*}[!t]
  \centering
  \begin{tikzpicture}[
  kairo/base, x=1cm, y=1cm,
  receive/.style={circle, minimum size=4.3pt, inner sep=0pt, line width=0.75pt, fill=white},
  resolve/.style={rectangle, minimum size=4.1pt, inner sep=0pt},
  same team/.style={draw=KairoBlue, line width=0.85pt, line cap=round},
  handover/.style={draw=KairoAmber, line width=0.85pt, dash pattern=on 3pt off 1.5pt,
    -{Stealth[length=4pt]}, shorten >=3pt}
]
  \path[use as bounding box] (0,0) rectangle (13.08,-2.75);
  % Use the same local time positions as the inter-case diagram.
  \foreach \windowid/\leftx/\windowtitle/\focalend/\focaly in {
    c/2.10/Earlier/2.84/-1.31,
    d/7.81/Later/4.49/-1.97
  } {
    \begin{scope}[shift={(\leftx,0)}]
      \node[font=\footnotesize\bfseries] at (2.635,-0.17)
        {\windowtitle\ window $W_{\mathrm{\windowid}}$};
      \path[fill=KairoBlueTint, rounded corners=2pt] (0,-1.58) rectangle (5.27,-0.50);
      \path[fill=KairoAmberTint, rounded corners=2pt] (0,-2.26) rectangle (5.27,-1.64);
      \coordinate (companion-start-\windowid) at (1.46,-0.81);
      \coordinate (companion-end-\windowid) at (3.56,-0.81);
      \coordinate (focal-start-\windowid) at (0.74,-1.31);
      \coordinate (focal-end-\windowid) at (\focalend,\focaly);
    \end{scope}
  }
  \node[anchor=east, font=\footnotesize\bfseries] at (1.96,-0.49) {Frontline};
  \node[anchor=east, text=KairoMuted] at (1.96,-0.81) {Companion};
  \node[anchor=east, text=KairoMuted] at (1.96,-1.31) {Focal case};
  \node[anchor=east, font=\footnotesize\bfseries] at (1.96,-1.97) {Specialist};
  \foreach \windowid in {c,d} {
    \draw[same team] (companion-start-\windowid) -- (companion-end-\windowid);
    \node[receive, draw=KairoBlue] at (companion-start-\windowid) {};
    \node[resolve, fill=KairoBlue] at (companion-end-\windowid) {};
  }
  \draw[same team] (focal-start-c) -- (focal-end-c);
  \draw[handover] (focal-start-d) -- (focal-end-d);
  \node[anchor=south, text=KairoAmber, inner sep=1pt] at (10.69,-1.59) {Handover};
  \node[receive, draw=KairoBlue] at (focal-start-c) {};
  \node[receive, draw=KairoBlue] at (focal-start-d) {};
  \node[resolve, fill=KairoBlue] at (focal-end-c) {};
  \node[resolve, fill=KairoAmber] at (focal-end-d) {};
  \node[receive, draw=KairoMuted] at (4.85,-2.60) {};
  \node[anchor=west, text=KairoMuted] at (4.98,-2.60) {Receive};
  \node[resolve, fill=KairoMuted] at (6.56,-2.60) {};
  \node[anchor=west, text=KairoMuted] at (6.69,-2.60) {Resolve};
  \draw[-{Stealth[length=4pt]}, draw=KairoMuted, line width=0.5pt]
    (8.19,-2.60) -- (8.74,-2.60);
  \node[anchor=west, text=KairoMuted] at (8.84,-2.60) {Time within window};
\end{tikzpicture}

  \caption{The focal case passes from Frontline to Specialist in $W_{\mathrm d}$. The handover is assigned to the window containing the Specialist event.}
  \label{fig:example-resource}
\end{figure*}

\begin{Example}[Inter-case features for a longer gap]\label{ex:inter-case}
Both $W_{\mathrm c}$ and $W_{\mathrm d}$ have the same numbers of arrivals, completions, and active cases. Each Resolve event contributes its gap from Receive to the mean for its window. The gaps are both $g$ in $W_{\mathrm c}$. In $W_{\mathrm d}$, the focal case contributes $G$ and the companion still contributes $g$, so the mean rises. \autoref{tab:example-inter-case} shows this change in timing alongside the matching counts. \autoref{fig:example-inter-case} shows which gaps enter each window's mean.
\end{Example}

\begin{table*}[!t]
  \caption{Selected inter-case features in consecutive feature windows. The later window has the same counts and a longer mean gap.}
  \label{tab:example-inter-case}
  \centering
  \small
  \begin{tabularx}{\linewidth}{@{}X*{2}{>{\centering\arraybackslash}p{0.17\linewidth}}@{}}
    \toprule
    \textbf{Feature} & $\boldsymbol{W_{\mathrm c}}$ & $\boldsymbol{W_{\mathrm d}}$ \\
    \midrule
    Case arrivals & Two & Two \\
    Case completions & Two & Two \\
    Active cases & Two & Two \\
    Events per active case & Two & Two \\
    Mean within-case event gap & $g$ & $\operatorname{mean}(g,G)$ \\
    \bottomrule
  \end{tabularx}
\end{table*}

\begin{figure*}[!t]
  \centering
  \begin{tikzpicture}[
  kairo/base, x=1cm, y=1cm,
  receive/.style={circle, minimum size=4.3pt, inner sep=0pt, draw=KairoTeal,
    line width=0.75pt, fill=white},
  resolve/.style={rectangle, minimum size=4.1pt, inner sep=0pt, fill=KairoTeal},
  gap/.style={draw=KairoTeal, line width=0.9pt, line cap=round, text=KairoTeal}
]
  \path[use as bounding box] (0,0) rectangle (13.08,-2.20);
  % Use the same local time positions as the resource diagram.
  \foreach \windowid/\leftx/\windowtitle/\focalend/\focalgap/\windowmean in {
    c/2.10/Earlier/2.84/g/g,
    d/7.81/Later/4.49/G/{\operatorname{mean}(g,G)}
  } {
    \begin{scope}[shift={(\leftx,0)}]
      \node[font=\footnotesize\bfseries] at (2.635,-0.17)
        {\windowtitle\ window $W_{\mathrm{\windowid}}$};
      \path[fill=KairoPaper, rounded corners=2pt] (0,-1.66) rectangle (5.27,-0.50);
      \draw[gap] (1.46,-0.91) -- node[midway, above=1pt, inner sep=1pt] {$g$} (3.56,-0.91);
      \draw[gap] (0.74,-1.43) -- node[midway, above=1pt, inner sep=1pt] {$\focalgap$}
        (\focalend,-1.43);
      \node[receive] at (1.46,-0.91) {};
      \node[resolve] at (3.56,-0.91) {};
      \node[receive] at (0.74,-1.43) {};
      \node[resolve] at (\focalend,-1.43) {};
      \path[fill=KairoTealTint, rounded corners=2pt] (0,-2.15) rectangle (5.27,-1.77);
      \node[text=KairoTeal] at (2.635,-1.96) {$\windowmean$};
    \end{scope}
  }
  \node[anchor=east] at (1.96,-0.91) {Companion};
  \node[anchor=east] at (1.96,-1.43) {Focal case};
  \node[anchor=east, font=\footnotesize\bfseries] at (1.96,-1.96) {Window mean};
\end{tikzpicture}

  \caption{The focal case's longer gap raises the window mean in $W_{\mathrm d}$. Open circles mark arrivals and filled squares mark completions.}
  \label{fig:example-inter-case}
\end{figure*}
